


\chapter{模板使用说明}


\section{表格}

\textbf{写在最开始}：由于一些实现上的问题，对于\textbf{确定非置底排版}的任何表格，用户在通过\shad{table}环境的选项指定可选择的排版模式时，\textbf{\color{DarkRed}不可提供\shad{b}模式}，否则表格的上间距将过窄；反之，对于页内\textbf{确定置底排版}的第一个表格，用户需要\textbf{\color{DarkRed}显式在选项中指定\shad{b}或\shad{!b}}，否则此表格上间距将过宽。

若出现意料之外的情况，用户可以通过在\shad{table}环境开始后和结束前的位置插入\shad{$\backslash$vspace*\{<距离长度>\}}来调整其上下间距，使之看起来协调。

\subsection{普通表格}

普通表格的排版本身无需多言，使用\shad{table}+\shad{tabular}环境即可，但是要注意三线表中的三条线分别需要使用\shad{$\backslash$toprule}、\shad{$\backslash$midrule}、\shad{$\backslash$bottomrule}生成，这样才符合研究生规范中对线粗的要求（\shad{1.5}磅、\shad{0.75}磅、\shad{1.5}磅）。注意不要用\shad{$\backslash$hline}。而对于本科生，学士学位论文规范要求表格中的线粗统一为\shad{0.5}磅。因此，在\shad{bachelor}和\shad{doublebachelor}选项下，\shad{$\backslash$toprule}、\shad{$\backslash$midrule}、\shad{$\backslash$bottomrule}和\shad{$\backslash$cmidrule}的线粗均设置为\shad{0.5}磅（2025.05.01调整）。

\begin{table}[htp]
    \caption{$\backslash$cmidrule示例}\label{tab: cmidrule示例}
    \begin{threeparttable}
        \begin{tabular}{ccccc}
            \toprule
            \multirow{2}{*}{Column0} &  \multicolumn{2}{c}{Column1\tnote{1}} & \multicolumn{2}{c}{Column2\tnote{2}} \\
            \cmidrule(lr){2-3}\cmidrule[2.5bp](l){4-5}
            ~     & subcolumn1 & subcolumn2 & subcolumn1 & subcolumn2 \\
            \midrule
            Row1  & element11 & element12 &element13 & element14 \\
            Row2  & element21 & element22 &element23 & element24 \\
            \cmidrule{2-3}\cmidrule[2.5bp]{4-5}
            Row3  & element31 & element32 &element33 & element34 \\
            \bottomrule
        \end{tabular}
        \begin{tablenotes}
            \item[1] 在\shad{bachelor}和\shad{doublebachelor}选项下，\shadcmd{cmidrule}默认线粗设置为0.5bp，而在其他论文类型下，默认值为0.75bp，两者规范的要求不同。可用第二项可选参数同时修剪掉框线的左右端点
            \item[2] 通过指定\shadcmd{cmidrule}的第一项可选参数调整线粗，并用第二项可选参数仅修剪掉框线的左端点
        \end{tablenotes}
    \end{threeparttable}
\end{table}

在需要为表格中的某些单元格添加水平框线时，应使用\newline\shadcmd{cmidrule[<线粗>](<修剪>)\{<起始列-终止列>\}}而非\shadcmd{cline\{<起始列-终止列>\}}。后者似乎无法调整线粗，也无法对框线的端点进行修剪。前者的第一项可选参数允许用户设置框线粗细，其默认值在\shad{bachelor}和\shad{doublebachelor}选项下设置为了\shad{0.5}磅，而在其他论文类型下设置为了\shad{0.75}磅。若非必要，用户无需设置该可选参数；第二项可选参数允许用户对框线的端点进行修剪，该选项可防止同行独立的相邻框线在视觉上连通到一起，参见表\ref{tab: cmidrule示例}中的示例。有关第二项可选参数可取的值，建议用户查阅\href{https://mirrors.sustech.edu.cn/CTAN/macros/latex/contrib/booktabs/booktabs.pdf}{\ttfamily\color{DarkRed}\CJKunderline*[thickness=0.5bp, format=\color{DarkRed}]{booktabs}}宏包的官方文档。

若要调整表格中单元格的水平对齐方式，自适应宽度的表格可以直接使用\shad{tabular}环境的\shad{l/c/r}参数；如果同时需要限定单元格宽度，则可使用\shad{p\{<宽度值>\}}、\shad{p\{<宽度值>\}<\{$\backslash$centering\}}、\shad{p\{<宽度值>\}<\{$\backslash$raggedleft\}}参数，效果见表\ref{tab: 单元格三种水平对齐方式}。

\begin{table}[!ht]
    \caption{\shad{单元格三种水平对齐方式}} \label{tab: 单元格三种水平对齐方式}
    \begin{tabular}{p{2cm} p{3cm}<{\centering} p{7cm}<{\raggedleft}}
        \toprule
        \textbf{姓名} & \textbf{所属势力} & \textbf{武功绝学} \\
        \midrule
        郭靖 & 重阳宫 & 降龙十八掌 \\
        黄蓉 & 丐帮 & 打狗棒法 \\
        洪七公 & 丐帮 & 降龙十八掌、打狗棒法 \\
        黄老邪 & 桃花岛 & 弹指神通、落英神剑掌、玉箫剑法 \\
        老顽童 & 重阳宫 & 左右互博术 \\
        一灯 & 云南大理 & 一阳指、千里传音 \\
        \bottomrule
    \end{tabular}
\end{table}

若要调整具有多行内容的固定列宽单元格的垂直对齐方式，可以使用\shad{tabular}环境的\shad{p\{<宽度值>\}}、\shad{m\{<宽度值>\}}、\shad{b\{<宽度值>\}}参数，\textbf{它们分别表示使用对应单元格内容的上线、中线、底线作为垂直对齐时的基准线}，\textcolor{DarkRed}{而非使单元格内容上端对齐、垂直居中对齐、下端对齐}，效果见表\ref{tab: 单元格三种垂直对齐方式}。如果这段说明不好理解，那可以参考这个：\href{https://tex.stackexchange.com/questions/35293}{\ttfamily\color{DarkRed}\CJKunderline*[thickness=0.5bp, format=\color{DarkRed}]{p, m and b columns in tables}}。

\begin{table}[!ht]
    \caption{\shad{单元格三种垂直对齐方式}} \label{tab: 单元格三种垂直对齐方式}
    \begin{tabular}{p{2cm} b{2cm} m{2cm}}
        \toprule
        \textbf{姓名} & \textbf{所属势力} & \textbf{武功绝学} \\
        \midrule
        洪七公 & 丐帮 & 降龙十八掌、打狗棒法 \\
        黄老邪黄老邪黄老邪 & 桃花岛桃花岛桃花岛 & 弹指神通、落英神剑掌、玉箫剑法 \\
        \bottomrule
    \end{tabular}
\end{table}

\newpage
\subsection{带附注表格}

更需要说明的是生成带附注的表格。本模板采用\shad{threeparttable}宏包实现将表格中的附注内容顶格排版在表格底部：
\begin{enumerate}
    \item 使用\shadcmd{tnote\{<label>\}}在表格中插入上标编号；
    \setcounter{enumi}{98}% 更改后续列表编号的起始值
    \item 使用\shad{tablenotes}环境在表格底部排版附注。该环境提供选项\shad{online}用于将附注文本前的标号从默认的上标样式（见表\ref{tab: 江湖势力背调}）更改为非上标样式（见表\ref{tab: 已习得武功}）。
\end{enumerate}

\begin{table}[!ht]
    \caption{江湖势力背调} \label{tab: 江湖势力背调}
    \begin{threeparttable}
        \begin{tabular}{p{2cm} p{3cm}<{\centering} p{7cm}<{\raggedleft}}
            \toprule
            \textbf{姓名} & \textbf{所属势力} & \textbf{武功绝学} \\
            \midrule
            郭靖 & 重阳宫 & 降龙十八掌 \\
            黄蓉 & 丐帮 & 打狗棒法 \\
            洪七公 & 丐帮 & 降龙十八掌、打狗棒法 \\
            黄老邪 & 桃花岛 & 弹指神通、落英神剑掌、玉箫剑法 \\
            老顽童 & 重阳宫 & 左右互博术\tnote{1} \\
            一灯 & 云南大理 & 一阳指\tnote{2}、千里传音 \\
            \bottomrule
        \end{tabular}
        \begin{tablenotes}
            \item[1] 左右互搏术是金庸小说《射雕英雄传》中「老顽童」周伯通在桃花岛的地洞中创出的武功，本质是一心二用，能够两手同时做不同的事情，在金庸武侠体系中是一门非常精妙的武学，其对于人物本身的战斗力加成堪称台阶性。
            \item[2] 云南大理段氏嫡传的武功，在点穴功夫中位居天下第一，运功后以右手食指点穴，出指可缓可快，缓时潇洒飘逸，快则疾如闪电，但着指之处，分毫不差。当与敌挣搏凶险之际，用此指法既可贴近径点敌人穴道，也可从远处欺近身去，一中即离，一攻而退，实为克敌保身的无上妙术。
        \end{tablenotes}
    \end{threeparttable}
\end{table}

\begin{table}[!ht]
    \caption{已习得武功} \label{tab: 已习得武功}
    \begin{threeparttable}
        \begin{tabular}{p{3cm} p{3cm} p{5cm}}
            \toprule
            \textbf{武功绝学} & \textbf{传授者} & \textbf{传授地点} \\
            \midrule
            蛤蟆功 & 欧阳锋 & 重阳山脉 \\
            九阴真经 & 小龙女 & 活死人墓 \\
            打狗棒法 & 洪七公、黄蓉 & 华山之巅、英雄大会 \\
            玉箫剑法 & 黄老邪 & 深山老林 \\
            黯然销魂掌\tnote{1} & 自创 & 海边 \\
            \bottomrule
        \end{tabular}
        \begin{tablenotes}[online]
            \item[1] 黯然销魂掌，是在杨过与小龙女离别后，认为今生再也见不到小龙女，悲从中来，由此创作了黯然销魂掌。黯然销魂掌和心情有关，此后杨过与小龙女重逢后，其心理愉悦，故使不出黯然销魂掌。
        \end{tablenotes}
    \end{threeparttable}
\end{table}

上述方式排版的带附注表格无法通过点击表格中的编号跳转到对应附注。为此，本模板提供命令\shadcmd{puttablenotelabel\{<标签>\}}和\shadcmd{tablenoteref\{<标签>\}}来实现该操作\textcolor{red}{（2025.01.05新增）}：

\begin{itemize}
    \item 首先将\shad{tablenotes}环境中手动设置的编号替换为\shadcmd{puttablenotelabel\{<标签>\}}；
    \item 然后在表格内容的对应位置使用\shadcmd{tablenoteref\{<标签>\}}即可。编号将自动生成，并按照使用\shadcmd{puttablenotelabel}的顺序递加。
\end{itemize}

这种方式由于需要建立交叉引用，通常需要用户编译两次。切记，\textbf{标签必须全文唯一}。表\ref{tab: 江湖势力背调（基于puttablenotelabel和tablenoteref）}提供了使用示例。

\begin{table}[!ht]
    \caption{江湖势力背调（基于\shadcmd{puttablenotelabel}和\shadcmd{tablenoteref}）} \label{tab: 江湖势力背调（基于puttablenotelabel和tablenoteref）}
    \begin{threeparttable}
        \begin{tabular}{p{2cm} p{3cm} p{7cm}}
            \toprule
            \textbf{姓名} & \textbf{所属势力} & \textbf{武功绝学} \\
            \midrule
            郭靖 & 重阳宫 & 降龙十八掌 \\
            黄蓉 & 丐帮 & 打狗棒法 \\
            洪七公 & 丐帮 & 降龙十八掌、打狗棒法 \\
            黄老邪 & 桃花岛 & 弹指神通、落英神剑掌、玉箫剑法 \\
            老顽童 & 重阳宫 & 左右互博术\tablenoteref{tn: 左右互搏术} \\
            一灯 & 云南大理 & 一阳指\tablenoteref{tn: 一阳指}、千里传音 \\
            \bottomrule
        \end{tabular}
        \begin{tablenotes}
            \item[\puttablenotelabel{tn: 左右互搏术}] 左右互搏术是金庸小说《射雕英雄传》中「老顽童」周伯通在桃花岛的地洞中创出的武功，本质是一心二用，能够两手同时做不同的事情，在金庸武侠体系中是一门非常精妙的武学，其对于人物本身的战斗力加成堪称台阶性。
            \item[\puttablenotelabel{tn: 一阳指}] 云南大理段氏嫡传的武功，在点穴功夫中位居天下第一，运功后以右手食指点穴，出指可缓可快，缓时潇洒飘逸，快则疾如闪电，但着指之处，分毫不差。当与敌挣搏凶险之际，用此指法既可贴近径点敌人穴道，也可从远处欺近身去，一中即离，一攻而退，实为克敌保身的无上妙术。
        \end{tablenotes}
    \end{threeparttable}
\end{table}


\clearpage
\subsection{跨页表格}

原则上，长度不足一页的表格不应跨页。而对于本身超过一页的表格，本模板使用\href{https://mirrors.tuna.tsinghua.edu.cn/CTAN/macros/latex/required/tools/longtable.pdf}{\ttfamily\color{DarkRed}\CJKunderline*[thickness=0.5bp, format=\color{DarkRed}]{longtable}}宏包提供的\shad{longtable}环境实现。用户需要了解\shad{longtable}环境的基本使用方法，它与\shad{tabular}环境的最大区别在于需要用户自行定义分页后的表题、表头以及表尾。

本模板提供了命令\shadcmd{CPcaption\{<当前表格总列数>\}\{<跨页表题>\}}来正确排版跨页之后的表题\textcolor{red}{（2025.01.05更新命令使用方式）}。\textbf{务必使用此命令}，否则跨页后的表题将会与表格内容采用相同的行距和段前段后，而非与规范中要求的表题格式保持一致；并且在跨页表题长度超过表宽时，无法产生预期的排版结果。

此外，\textbf{\textcolor{DarkRed}{不应将\shad{longtable}环境嵌套在\shad{table}等浮动环境中}}，否则长表格将无法正常跨页。具体细节参见本小节示例表\ref{tab: 中国计算机学会部分推荐期刊及会议}和表\ref{tab: 中国计算机学会部分推荐期刊及会议简表}。

% \newpage

\begin{longtable}{p{2em} p{4.5em} p{11em} p{6em} p{11em}}
    \caption{中国计算机学会部分推荐期刊及会议} \label{tab: 中国计算机学会部分推荐期刊及会议} \\
    
    \toprule
    \textbf{序号} & \textbf{刊物简称} & \textbf{刊物全称} & \textbf{出版社} & \textbf{网址} \\
    \midrule
    \endfirsthead
    
    % 在这里设计首页以外的表题和表头
    \CPcaption{5}{中国计算机学会部分推荐期刊及会议}\\
    \toprule
    \textbf{序号} & \textbf{刊物简称} & \textbf{刊物全称} & \textbf{出版社} & \textbf{网址} \\
    \midrule
    \endhead
    
    % 在这里设计首页以外的表尾
    \bottomrule
    \multicolumn{5}{l}{续下页} \\  % 如不希望跨页表尾显示任何内容则注释掉即可
    \endfoot
    
    \bottomrule
    \endlastfoot
    
    1 & JSAC & IEEE Journal on Selected Areas in Communications & IEEE & http://dblp.uni-trier.de/db/journals/jsac/ \\
    2 & TMC & IEEE Transactions on Mobile Computing & IEEE & http://dblp.uni-trier.de/db/journals/tmc/ \\
    3 & TON & IEEE/ACM Transactions on Networking & IEEE/ACM & http://dblp.uni-trier.de/db/journals/ton/ \\
    1 & TOIT & ACM Transactions on Internet Technology & ACM & http://dblp.uni-trier.de/db/journals/toit/ \\
    2 & TOMM & ACM Transactions on Multimedia Computing, Communications and Applications & ACM & http://dblp.uni-trier.de/db/journals/tomccap/ \\
    3 & TOSN & ACM Transactions on Sensor Networks & ACM & http://dblp.uni-trier.de/db/journals/tosn/ \\
    4 & CN & Computer Networks & Elsevier & http://dblp.uni-trier.de/db/journals/cn/ \\
    5 & TCOM & IEEE Transactions on Communications & IEEE & http://dblp.uni-trier.de/db/journals/tcom/ \\
    6 & TWC & IEEE Transactions on Wireless Communications & IEEE & http://dblp.uni-trier.de/db/journals/twc/ \\
    1 & & Ad Hoc Networks & Elsevier & http://dblp.uni-trier.de/db/journals/adhoc/ \\
    2 & CC & Computer Communications & Elsevier & http://dblp.uni-trier.de/db/journals/comcom/ \\
    3 & TNSM & IEEE Transactions on Network and Service Management & IEEE & http://dblp.uni-trier.de/db/journals/tnsm/ \\
    4 & & IET Communications & IET & http://dblp.uni-trier.de/db/journals/iet-com/ \\
    5 & JNCA & Journal of Network and Computer Applications & Elsevier & http://dblp.uni-trier.de/db/journals/jnca/ \\
    6 & MONET & Mobile Networks and Applications & Springer & http://dblp.uni-trier.de/db/journals/monet/ \\
    7 & & Networks & Wiley & http://dblp.uni-trier.de/db/journals/networks/ \\
    8 & PPNA & Peer-to-Peer Networking and Applications & Springer & http://dblp.uni-trier.de/db/journals/ppna/ \\
    9 & WCMC & Wireless Communications and Mobile Computing & Wiley & http://dblp.uni-trier.de/db/journals/wicomm/ \\
    10 & & Wireless Networks & Springer & http://dblp.uni-trier.de/db/journals/winet/ \\
    11 & IOT & IEEE Internet of Things Journal & IEEE & https://dblp.org/db/journals/ iotj/index.html \\
    1 & SIGCOMM & ACM International Conference on Applications, Technologies, Architectures, and Protocols for Computer Communication & ACM & http://dblp.uni-trier.de/db/conf/sigcomm/ index.html \\
    2 & MobiCom & ACM International Conference on Mobile Computing and Networking & ACM & http://dblp.uni-trier.de/db/conf/mobicom/ \\
    3 & INFOCOM & IEEE International Conference on Computer Communications & IEEE & http://dblp.uni-trier.de/db/conf/infocom/ \\
    4 & NSDI & Symposium on Network System Design and Implementation & USENIX & http://dblp.uni-trier.de/db/conf/nsdi/ \\
    1 & SenSys & ACM Conference on Embedded Networked Sensor Systems & ACM & http://dblp.uni-trier.de/db/conf/sensys/ \\
    2 & CoNEXT & ACM International Conference on Emerging Networking Experiments and Technologies & ACM & http://dblp.uni-trier.de/db/conf/conext/ \\
    3 & SECON & IEEE International Conference on Sensing, Communication, and Networking & IEEE & http://dblp.uni-trier.de/db/conf/secon/ \\
    4 & IPSN & International Conference on Information Processing in Sensor Networks & IEEE/ACM & http://dblp.uni-trier.de/db/conf/ipsn/ \\
    5 & MobiSys & ACM International Conference on Mobile Systems, Applications, and Services & ACM & http://dblp.uni-trier.de/db/conf/mobisys/ \\
    6 & ICNP & IEEE International Conference on Network Protocols & IEEE & http://dblp.uni-trier.de/db/conf/icnp/ \\
    7 & MobiHoc & International Symposium on Theory, Algorithmic Foundations, and Protocol Design for Mobile Networks and Mobile Computing & ACM/IEEE & http://dblp.uni-trier.de/db/conf/mobihoc/ \\
    8 & NOSSDAV & International Workshop on Network and Operating System Support for Digital Audio and Video & ACM & http://dblp.uni-trier.de/db/conf/nossdav/ \\
    9 & IWQoS & IEEE/ACM International Workshop on Quality of Service & IEEE & http://dblp.uni-trier.de/db/conf/iwqos/ \\
    10 & IMC & ACM Internet Measurement Conference & ACM/USENIX & http://dblp.uni-trier.de/db/conf/imc/ \\
    
    
\end{longtable}

\newpage

\begin{longtable}{p{2em} p{4.5em}}
    \caption{中国计算机学会部分推荐期刊及会议简表（用于测试跨页表宽低于表题长度的情况）} \label{tab: 中国计算机学会部分推荐期刊及会议简表} \\
    
    \toprule
    \textbf{序号} & \textbf{刊物简称} \\
    \midrule
    \endfirsthead
    
    % 在这里设计首页以外的表题和表头
    \CPcaption{2}{中国计算机学会部分推荐期刊及会议简表（用于测试跨页表宽低于表题长度的情况）}\\
    \toprule
    \textbf{序号} & \textbf{刊物简称} \\
    \midrule
    \endhead
    
    % 在这里设计首页以外的表尾
    \bottomrule
    \multicolumn{2}{r}{续下页} \\  % 如不希望跨页表尾显示任何内容则注释掉即可
    \endfoot
    
    \bottomrule
    \endlastfoot
    
    1 & JSAC \\
    2 & TMC \\
    3 & TON \\
    1 & TOIT \\
    2 & TOMM \\
    3 & TOSN \\
    4 & CN \\
    5 & TCOM \\
    6 & TWC \\
    2 & CC \\
    3 & TNSM \\
    5 & JNCA \\
    6 & MONET \\
    8 & PPNA \\
    9 & WCMC \\
    11 & IOT \\
    1 & SIGCOMM \\
    2 & MobiCom \\
    3 & INFOCOM \\
    4 & NSDI \\
    1 & SenSys \\
    2 & CoNEXT \\
    3 & SECON \\
    4 & IPSN \\
    5 & MobiSys \\
    6 & ICNP \\
    7 & MobiHoc \\
    8 & NOSSDAV \\
    9 & IWQoS \\
    10 & IMC \\
    1 & JSAC \\
    2 & TMC \\
    3 & TON \\
    1 & TOIT \\
    2 & TOMM \\
    3 & TOSN \\
    4 & CN \\
    5 & TCOM \\
    6 & TWC \\
    2 & CC \\
    3 & TNSM \\
    5 & JNCA \\
    6 & MONET \\
    8 & PPNA \\
    9 & WCMC \\
    11 & IOT \\
    1 & SIGCOMM \\
    2 & MobiCom \\
    3 & INFOCOM \\
    4 & NSDI \\
    1 & SenSys \\
    2 & CoNEXT \\
    3 & SECON \\
    4 & IPSN \\
    5 & MobiSys \\
    6 & ICNP \\
    7 & MobiHoc \\
    8 & NOSSDAV \\
    9 & IWQoS \\
    10 & IMC \\
\end{longtable}

\clearpage
\subsection{跨页带附注表格（2025.01.05）}

很遗憾，\shad{threeparttable}无法做到跨页，而\shad{longtable}又无法像前者那样稳定完美地排版附注。无奈之下，本模板提供了一种繁琐但可行的做法。

思路是利用\shad{longtable}提供的表尾自定义功能来插入附注，即用户需要修改\shad{longtable}在\shadcmd{endlastfoot}前对表格尾部的设置。为此，模板提供命令：

\shadcmd{tablenotetext(online)(<附注编号悬挂距离>)[<附注上方垂直间隔>]\{<附注总宽度>\}\{<附注标签>\}\{<附注内容>\}}。

\begin{itemize}
    \item 此命令的第一项可选参数以\shad{()}标识，它仅接受参数\shad{online}，作用是将附注编号从默认的上标形式更改为行内形式，即模仿\shad{tablenotes}环境的选项。此间差异可参考表\ref{tab: 跨页带附注表格示例}的附注；
    \item 第二项可选参数仅在设置了第一项可选参数为\shad{online}时有效，用于调整附注编号的悬挂缩进距离，默认是\shad{1em}。当表格的附注数量过多时，可通过手动调整该可选参数来避免过长的编号与后方文字重叠，见表\ref{tab: 跨页带附注表格恰巧在附注开始处换页示例}；
    \item 第三项可选参数以\shad{[]}标识，在生成的附注与上方内容间的垂直距离不合适时，可通过此可选参数手动对其进行调整，默认为\shad{0bp}；
    \item 第四项强制参数对应附注整体的宽度，其取值需要根据表格排版后的宽度反复调整，直至附注与表格尾线等宽为止。这的确很繁琐，但我实在能力有限，想不出更好的办法，自动确定表格的实际宽度真的很难；
    \item 第五项强制参数负责设置附注编号对应的标签，以便后续在表格中使用\shadcmd{tablenoteref\{\}}进行引用，标签同样需要全文唯一；
    \item 第六项强制参数对应附注的实际内容。
\end{itemize}

\textcolor{red}{\textbf{注意：}}使用\shadcmd{tablenotetext}命令的前提是对表格首列进行特定设置，用户\textbf{必须}通过\shad{p\{<长度>\}}或\shad{m\{<长度>\}}来人为指定\shad{longtable}的首列宽度，\textcolor{red}{\textbf{切不可}}将之设置为\shad{c}、\shad{r}或\shad{l}。另外，表格中最后一条\shadcmd{tablenotetext}之后不需要\shadcmd{\shadcmd{}}，否则表尾与下方文本的间隔将偏大。

可能出现一种极端情况：跨页表格恰好在附注开始处发生了分页，此时也会产生另类的排版结果。要解决该问题，用户得手动在表格内容的适当位置使用\href{https://mirrors.tuna.tsinghua.edu.cn/CTAN/macros/latex/required/tools/longtable.pdf}{\ttfamily\color{DarkRed}\CJKunderline*[thickness=0.5bp, format=\color{DarkRed}]{longtable}}宏包提供的\shadcmd{pagebreak}命令提前断页，代价是前一页底部可能有更多空白，参考表\ref{tab: 跨页带附注表格恰巧在附注开始处换页示例}中的做法。如果你运气尤其差，附注特别长，而且在附注中间跨页了，那我实在无能为力了。

必须要承认，\shadcmd{tablenotetext}命令很不稳定。其参数在不同的表格中需要特调，甚至在同一表格的不同排版设置下都是如此，可谓一表一参。如遇到不得不用的时候，用户必须要投入很多精力。可惜，这已经是我目前所能做到的极限。


\newpage

\begin{longtable}{p{2em} p{4.5em} p{20em} p{6em}}
    \caption{跨页带附注表格示例} \label{tab: 跨页带附注表格示例} \\
    
    \toprule
    \textbf{序号} & \textbf{刊物简称} & \textbf{刊物全称} & \textbf{出版社} \\
    \midrule
    \endfirsthead
    
    % 在这里设计首页以外的表题和表头
    \CPcaption{4}{跨页带附注表格示例}\\
    \toprule
    \textbf{序号} & \textbf{刊物简称} & \textbf{刊物全称} & \textbf{出版社} \\
    \midrule
    \endhead
    
    % 在这里设计首页以外的表尾
    \bottomrule
    \multicolumn{4}{l}{续下页} \\  % 如不希望跨页表尾显示任何内容则注释掉即可
    \endfoot
    
    \bottomrule
    \tablenotetext[-7bp]{37.2em}{tn: 手动跨页表格附注标签IEEE}{IEEE是指电气和电子工程师学会，是一个国际性的专业学会，以促进电气工程、电子工程、计算机科学和相关领域的科学和技术发展为宗旨。成立于1884年，总部位于美国纽约。IEEE 的会员包括来自世界各地的专业人士、工程师、学者和学生，是全球最大的技术专业组织之一。} \\
    \tablenotetext(online)[6bp]{37.2em}{tn: 手动跨页表格附注标签ACM}{ACM是指国际计算机学会，成立于1947年，是一个国际性的科技教育组织，是世界上第一个科学性及教育性计算机学会，总部设在美国纽约。国际计算机学会是世界上最大的计算机领域专业性学术组织，汇集了国际计算机领域教育家，研究人员，工业界人士及学生。ACM致力于提高在中国的活动的规格与影响力。在此基础上，学会成立了ACM中国理事会，为在中国的学会会员与学会活动提供支持。}% 注意这里不需要\\
    \endlastfoot
    
    1 & JSAC & IEEE Journal on Selected Areas in Communications & IEEE \\
    2 & TMC & IEEE Transactions on Mobile Computing & IEEE \\
    3 & TON & IEEE/ACM Transactions on Networking & IEEE/ACM \\
    1 & TOIT & ACM Transactions on Internet Technology & ACM \\
    2 & TOMM & ACM Transactions on Multimedia Computing, Communications and Applications & ACM \\
    3 & TOSN & ACM Transactions on Sensor Networks & ACM \\
    4 & CN & Computer Networks & Elsevier \\
    5 & TCOM & IEEE Transactions on Communications & IEEE \\
    6 & TWC & IEEE Transactions on Wireless Communications & IEEE \\
    1 & & Ad Hoc Networks & Elsevier \\
    2 & CC & Computer Communications & Elsevier \\
    3 & TNSM & IEEE Transactions on Network and Service Management & IEEE \\
    4 & & IET Communications & IET \\
    5 & JNCA & Journal of Network and Computer Applications & Elsevier \\
    6 & MONET & Mobile Networks and Applications & Springer \\
    7 & & Networks & Wiley \\
    8 & PPNA & Peer-to-Peer Networking and Applications & Springer \\
    9 & WCMC & Wireless Communications and Mobile Computing & Wiley \\
    10 & & Wireless Networks & Springer \\
    11 & IOT & IEEE Internet of Things Journal & IEEE \\
    1 & SIGCOMM & ACM International Conference on Applications, Technologies, Architectures, and Protocols for Computer Communication & ACM \\
    2 & MobiCom & ACM International Conference on Mobile Computing and Networking & ACM \\
    3 & INFOCOM & IEEE International Conference on Computer Communications & IEEE \\
    4 & NSDI & Symposium on Network System Design and Implementation & USENIX \\
    1 & SenSys & ACM Conference on Embedded Networked Sensor Systems & ACM \\
    2 & CoNEXT & ACM International Conference on Emerging Networking Experiments and Technologies & ACM \\
    3 & SECON & IEEE International Conference on Sensing, Communication, and Networking & IEEE \\
    4 & IPSN & International Conference on Information Processing in Sensor Networks & IEEE/ACM \\
    5 & MobiSys & ACM International Conference on Mobile Systems, Applications, and Services & ACM \\
    6 & ICNP & IEEE International Conference on Network Protocols & IEEE\tablenoteref{tn: 手动跨页表格附注标签IEEE} \\
    7 & MobiHoc & International Symposium on Theory, Algorithmic Foundations, and Protocol Design for Mobile Networks and Mobile Computing & ACM/IEEE \\
    8 & NOSSDAV & International Workshop on Network and Operating System Support for Digital Audio and Video & ACM\tablenoteref{tn: 手动跨页表格附注标签ACM} \\
    9 & IWQoS & IEEE/ACM International Workshop on Quality of Service & IEEE \\
    10 & IMC & ACM Internet Measurement Conference & ACM/USENIX \\
    
    
\end{longtable}

表格后参照文本表格后参照文本表格后参照文本表格后参照文本表格后参照文本表格后参照文本表格后参照文本表格后参照文本表格后参照文本表格后参照文本表格后参照文本表格后参照文本表格后参照文本表格后参照文本表格后参照文本表格后参照文本表格后参照文本表格后参照文本表格后参照文本表格后参照文本表格后参照文本

\newpage

\begin{longtable}{m{2em}<{\centering} p{4.5em} p{15em} p{6em}}
    \caption{跨页带附注表格恰巧在附注开始处换页示例} \label{tab: 跨页带附注表格恰巧在附注开始处换页示例} \\
    
    \toprule
    \textbf{序号} & \textbf{刊物简称} & \textbf{刊物全称} & \textbf{出版社} \\
    \midrule
    \endfirsthead
    
    % 在这里设计首页以外的表题和表头
    \CPcaption{4}{跨页带附注表格恰巧在附注中需要换页示例}\\
    \toprule
    \textbf{序号} & \textbf{刊物简称} & \textbf{刊物全称} & \textbf{出版社} \\
    \midrule
    \endhead
    
    % 在这里设计首页以外的表尾
    \bottomrule
    \multicolumn{4}{l}{续下页} \\  % 如不希望跨页表尾显示任何内容则注释掉即可
    \endfoot
    
    \bottomrule
    \tablenotetext[5bp]{32.1em}{tn: 手动跨页表格附注标签Elsevier}{爱思唯尔，创办于1880年，属于RELX集团旗下，总部位于阿姆斯特丹。爱思唯尔是一家荷兰的国际化多媒体出版集团，主要为科学家、研究人员、学生、医学以及信息处理的专业人士提供信息产品和革新性工具。爱思唯尔是全球领先的科学与医学信息服务机构，旗下出版《柳叶刀》《细胞》等2800多种学术期刊。} \setcounter{tablenote}{99}\\
    \tablenotetext(online)(2em)[5bp]{32.1em}{tn: 手动跨页表格附注标签USENIX}{USENIX成立于1975年，当时的名字叫做Unix用户群。它的主要目的是学习及开发Unix以及类似系统。1977年六月，美国电话电报公司的律师告诉用户群他们不能继续使用UNIX这个名字，因为UNIX是美国电话电报公司所拥有的一个商标。所以这个用户群更名成USENIX.从那以后，USENIX逐渐发展成一个倍受尊敬的由计算机操作系统用户，开发者和研究者所组成的机构。}% 注意这里不需要\\
    \endlastfoot
    
    1 & JSAC & IEEE Journal on Selected Areas in Communications & IEEE \\
    2 & TMC & IEEE Transactions on Mobile Computing & IEEE \\
    3 & TON & IEEE/ACM Transactions on Networking & IEEE/ACM \\
    1 & TOIT & ACM Transactions on Internet Technology & ACM \\
    2 & TOMM & ACM Transactions on Multimedia Computing, Communications and Applications & ACM \\
    3 & TOSN & ACM Transactions on Sensor Networks & ACM \\
    4 & CN & Computer Networks & Elsevier \\
    5 & TCOM & IEEE Transactions on Communications & IEEE \\
    6 & TWC & IEEE Transactions on Wireless Communications & IEEE \\
    1 & & Ad Hoc Networks & Elsevier \\
    2 & CC & Computer Communications & Elsevier\tablenoteref{tn: 手动跨页表格附注标签Elsevier} \\
    3 & TNSM & IEEE Transactions on Network and Service Management & IEEE \\
    % 4 & & IET Communications & IET \\
    % 5 & JNCA & Journal of Network and Computer Applications & Elsevier \\
    % 6 & MONET & Mobile Networks and Applications & Springer \\
    % 7 & & Networks & Wiley \\
    % 8 & PPNA & Peer-to-Peer Networking and Applications & Springer \\
    % 9 & WCMC & Wireless Communications and Mobile Computing & Wiley \\
    % 10 & & Wireless Networks & Springer \\
    % 11 & IOT & IEEE Internet of Things Journal & IEEE \\
    1 & SIGCOMM & ACM International Conference on Applications, Technologies, Architectures, and Protocols for Computer Communication & ACM \\
    2 & MobiCom & ACM International Conference on Mobile Computing and Networking & ACM \\
    3 & INFOCOM & IEEE International Conference on Computer Communications & IEEE \\
    4 & NSDI & Symposium on Network System Design and Implementation & USENIX\tablenoteref{tn: 手动跨页表格附注标签USENIX} \\
    1 & SenSys & ACM Conference on Embedded Networked Sensor Systems & ACM \\
    \pagebreak  % 用户可以尝试注释掉这条命令看看现象
    2 & CoNEXT & ACM International Conference on Emerging Networking Experiments and Technologies & ACM \\
    % 3 & SECON & IEEE International Conference on Sensing, Communication, and Networking & IEEE \\
    % 4 & IPSN & International Conference on Information Processing in Sensor Networks & IEEE/ACM \\
    % 5 & MobiSys & ACM International Conference on Mobile Systems, Applications, and Services & ACM \\
    % 6 & ICNP & IEEE International Conference on Network Protocols & IEEE \\
    % 7 & MobiHoc & International Symposium on Theory, Algorithmic Foundations, and Protocol Design for Mobile Networks and Mobile Computing & ACM/IEEE \\
    % 8 & NOSSDAV & International Workshop on Network and Operating System Support for Digital Audio and Video & ACM \\
    % 9 & IWQoS & IEEE/ACM International Workshop on Quality of Service & IEEE \\
    % 10 & IMC & ACM Internet Measurement Conference & ACM/USENIX \\
    
    
    
\end{longtable}

表格后参照文本表格后参照文本表格后参照文本表格后参照文本表格后参照文本表格后参照文本表格后参照文本表格后参照文本表格后参照文本表格后参照文本表格后参照文本表格后参照文本表格后参照文本表格后参照文本表格后参照文本表格后参照文本表格后参照文本表格后参照文本表格后参照文本表格后参照文本表格后参照文本

\clearpage

\section{伪代码}

除了\href{https://mirrors.sustech.edu.cn/CTAN/macros/latex/contrib/algorithm2e/doc/algorithm2e.pdf}{\ttfamily\color{DarkRed}\CJKunderline*[thickness=0.5bp, format=\color{DarkRed}]{algorithm2e}}宏包本身提供的各种条件、循环语句，本模板基于宏包提供的接口，追加了\shad{Do While}和\shad{Loop}循环语句：
\begin{itemize}
    \item \shadcmd{DoWhile(<紧跟关键字do的文本，可用于添加注释>)\{<循环条件>\}\{<循环体>\}}
    \item \shadcmd{Loop(<紧跟关键字loop的文本，可用于添加注释>)\{<循环体>\}}
\end{itemize}


此外，基于调整后的\shad{algorithm2e}环境，本模板进一步封装了\shad{algo}环境，它将生成比\shad{algorithm}环境更\textbf{\color{DarkRed}窄}的伪码浮动区域。除了接受浮动可选参数\shad{[htbp]}，\shad{algo}环境还支持另一可选参数\shad{(<伪码距正文文本边界的总距离>)}：该参数控制浮动体距正文文本边界的总距离，默认\shad{4em}，即单边缩进\shad{2em}，与其下首行文本对齐。两项可选参数可以单独或同时使用，\textbf{同时使用时的顺序必须与下方示例保持一致：}

\begin{verbatim}
    \begin{algo}[<浮动选项>](<伪码距正文文本边界的总距离>)
        .....
    \end{algo}
\end{verbatim}

算法\ref{alg: algorithm环境伪码示例}和算法\ref{alg: algo环境伪码示例}分别展示了两种环境默认生成的伪码样式；过程\ref{alg: algorithm环境修改伪码标签示例}和过程\ref{alg: algo环境修改伪码标签并调整宽度示例}展示了如何修改伪码中的一些标签，以及调整\shad{algo}伪码宽度的具体做法。


\begin{algorithm}[!h]
    \caption{algorithm环境伪码示例} \label{alg: algorithm环境伪码示例}
    \Input{1) 输入1;\newline 2) 输入2。}
    \Output{输出结果。}
    伪码行1。
    
    \For(\tcc*[f]{循环条件注释1}){循环条件1}{
        伪码行2。
        
        \tcp{注释2}
        伪码行3。
        
        \DoWhile(\tcc*[f]{循环条件注释3}){循环条件2}{
            伪码行4。
        }
        
        \tcc{loop循环}
        \Loop(\tcc*[f]{注释4}){
            循环体1。
        }
        
        \Repeat(\tcc*[f]{循环条件注释5}){循环条件3}{
            循环体2。
        }
        
        \tcp{if-elseif-else结构示例}
        \uIf(\tcc*[f]{条件注释6}){条件语句5}{
            条件语句5为真，伪码行5。
        }
        \uElseIf(\tcc*[f]{elseif条件语句}){条件语句6}{
            条件语句6为真，伪码行6。
        }
        \Else{
            条件5和6均为假，伪码行7。\tcp*[f]{else代码内容}
        }
        
        \If(\tcc*[f]{条件注释7}){条件语句7}{
            伪码行8。
        }
    }
    \textbf{return} 算法结果。
\end{algorithm}

\begin{algorithm}[!h]
    \renewcommand{\algorithmcfname}{过程}  % 修改伪码标签需要在\caption{}之前
    \caption{algorithm环境临时修改伪码标签示例} \label{alg: algorithm环境修改伪码标签示例}
    \SetKwInOut{Input}{In}
    \SetKwInOut{Output}{Out}
    \Input{1) 输入1; 2) 输入2。}
    \Output{输出结果。}
    伪码行1。
    
    \For(\tcc*[f]{循环条件注释1}){循环条件1}{
        伪码行2。
        
        \tcp{注释2}
        伪码行3。
        
        \DoWhile(\tcc*[f]{循环条件注释3}){循环条件2}{
            伪码行4。
        }
        
        \tcc{loop循环}
        \Loop(\tcc*[f]{注释4}){
            循环体1。
        }
        
        \Repeat(\tcc*[f]{循环条件注释5}){循环条件3}{
            循环体2。
        }
        \eIf(\tcc*[f]{条件注释6}){条件语句6}{
            为真，伪码行5。
        }{
            条件为假，伪码行6。\tcp*[f]{else代码内容}
        }
        
        \If(\tcc*[f]{条件注释7}){条件语句7}{
            伪码行7。
        }
    }
    \textbf{return} 算法结果。
\end{algorithm}

\begin{algo}[!h]
    \caption{algo环境伪码示例} \label{alg: algo环境伪码示例}
    \Input{1) 输入1;\newline 2) 输入2。}
    \Output{输出结果。}
    伪码行1。
    
    \For(\tcc*[f]{循环条件注释1}){循环条件1}{
        伪码行2。
        
        \tcp{注释2}
        伪码行3。
        
        \DoWhile(\tcc*[f]{循环条件注释3}){循环条件2}{
            伪码行4。
        }
        
        \tcc{loop循环}
        \Loop(\tcc*[f]{注释4}){
            循环体1。
        }
        
        \Repeat(\tcc*[f]{循环条件注释5}){循环条件3}{
            循环体2。
        }
        \eIf(\tcc*[f]{条件注释6}){条件语句6}{
            为真，伪码行5。
        }{
            条件为假，伪码行6。\tcp*[f]{else代码内容}
        }
        
        \If(\tcc*[f]{条件注释7}){条件语句7}{
            伪码行7。
        }
    }
    \textbf{return} 算法结果。
\end{algo}

\begin{algo}[!h](8em)
    \renewcommand{\algorithmcfname}{过程}  % 修改伪码标签需要在\caption{}之前
    \caption{algo环境临时修改伪码标签并调整宽度示例} \label{alg: algo环境修改伪码标签并调整宽度示例}
    \SetKwInOut{Input}{In}
    \SetKwInOut{Output}{Out}
    \Input{1) 输入1; 2) 输入2。}
    \Output{输出结果。}
    伪码行1。
    
    \For(\tcc*[f]{循环条件注释1}){循环条件1}{
        伪码行2。
        
        \tcp{注释2}
        伪码行3。
        
        \DoWhile(\tcc*[f]{循环条件注释3}){循环条件2}{
            伪码行4。
        }
        
        \tcc{loop循环}
        \Loop(\tcc*[f]{注释4}){
            循环体1。
        }
        
        \Repeat(\tcc*[f]{循环条件注释5}){循环条件3}{
            循环体2。
        }
        \eIf(\tcc*[f]{条件注释6}){条件语句6}{
            为真，伪码行5。
        }{
            条件为假，伪码行6。\tcp*[f]{else代码内容}
        }
        
        \If(\tcc*[f]{条件注释7}){条件语句7}{
            伪码行7。
        }
    }
    \textbf{return} 算法结果。
\end{algo}

\clearpage
\section{定义、公理、定理、命题、推论、引理、示例、假设、注、证明}

本模板分别定义了环境：\shad{definition}、\shad{axiom}、\shad{theorem}、\shad{proposition}、\shad{corollary}、\shad{lemma}、\shad{example}、\shad{assumption}、\shad{annotation}和\shad{proof}。示例如下：

% 云南大理段氏嫡传的武功，在点穴功夫中位居天下第一，运功后以右手食指点穴，出指可缓可快，缓时潇洒飘逸，快则疾如闪电，但着指之处，分毫不差。当与敌挣搏凶险之际，用此指法既可贴近径点敌人穴道，也可从远处欺近身去，一中即离，一攻而退，实为克敌保身的无上妙术。



\begin{axiom}[具体名称]
    云南大理段氏嫡传的武功，在点穴功夫中位居天下第一，运功后以右手食指点穴，出指可缓可快，缓时潇洒飘逸，快则疾如闪电。
\end{axiom}

\begin{theorem}[具体名称]
    云南大理段氏嫡传的武功，在点穴功夫中位居天下第一，运功后以右手食指点穴，出指可缓可快，缓时潇洒飘逸，快则疾如闪电。

    \begin{enumerate}
        \item 当与敌挣搏凶险之际，用此指法既可贴近径点敌人穴道，也可从远处欺近身去，一中即离，一攻而退，实为克敌保身的无上妙术。
        \begin{enumerate}
            \item 当与敌挣搏凶险之际，用此指法既可贴近径点敌人穴道，也可从远处欺近身去，一中即离，一攻而退，实为克敌保身的无上妙术。
        \end{enumerate}
    \end{enumerate}
\end{theorem}

\begin{proposition}[具体名称]
    云南大理段氏嫡传的武功，在点穴功夫中位居天下第一，运功后以右手食指点穴，出指可缓可快，缓时潇洒飘逸，快则疾如闪电。
\end{proposition}

\begin{corollary}[具体名称]
    云南大理段氏嫡传的武功，在点穴功夫中位居天下第一，运功后以右手食指点穴，出指可缓可快，缓时潇洒飘逸，快则疾如闪电。
\end{corollary}

\begin{lemma}[具体名称]
    云南大理段氏嫡传的武功，在点穴功夫中位居天下第一，运功后以右手食指点穴，出指可缓可快，缓时潇洒飘逸，快则疾如闪电。
\end{lemma}

\begin{example}[具体名称]
    云南大理段氏嫡传的武功，在点穴功夫中位居天下第一，运功后以右手食指点穴，出指可缓可快，缓时潇洒飘逸，快则疾如闪电。
\end{example}

\begin{assumption}[具体名称]
    云南大理段氏嫡传的武功，在点穴功夫中位居天下第一，运功后以右手食指点穴，出指可缓可快，缓时潇洒飘逸，快则疾如闪电。
\end{assumption}

\begin{annotation}[具体名称]
    云南大理段氏嫡传的武功，在点穴功夫中位居天下第一，运功后以右手食指点穴，出指可缓可快，缓时潇洒飘逸，快则疾如闪电。
\end{annotation}

\begin{proof}
    云南大理段氏嫡传的武功，在点穴功夫中位居天下第一，运功后以右手食指点穴，出指可缓可快，缓时潇洒飘逸，快则疾如闪电。
\end{proof}

\newpage

\section{脚注}

本模板使用包含了带圈数字的字体来替换LaTeX绘制的带圈数字，提供了充足的带圈编号数量，同时保证了带圈脚注编号足够优雅。

在正文中加入脚注直接在需要放置脚注标签的位置使用\shadcmd{footnote\{<脚注内容>\}}即可。

在其他环境中，如表格，则需要需要使用\shadcmd{footnotemark}配合\shadcmd{footnotetext\{<脚注文本>\}}。在需要放置脚注标签的位置使用\shadcmd{footnotemark}，然后在环境外使用\shadcmd{footnotetext\{<脚注文本>\}}指明脚注内容\footnote{更详细的使用方法参考\href{https://blog.csdn.net/xovee/article/details/127563209}{\ttfamily\color{DarkRed}\CJKunderline*[thickness=0.5bp, format=\color{DarkRed}]{LaTeX脚注}}。冗余文本用于展示脚注内容发生换行后的情况；冗余文本用于展示脚注内容发生换行后的情况；冗余文本用于展示脚注内容发生换行后的情况。}。


\section{模板中的各种编号}

“标题”、“图片”、“表格”、“伪码”、“公式”、“定义”、“定理”、“命题”、“推论”、“引理”、“证明”、“脚注”这些文档元素均可自行计算并生成编号，无需使用者费心。\textbf{但形如\textcolor{DodgerBlue}{(1-1a)}的子公式编号不能完全自动生成}，为此，模板提供了\shadcmd{subeqtag[<子公式编号标签>]}命令。

在为数学模型的约束创建编号时，常见的方式可能是使用\shadcmd{tag\{\}}命令直接指定编号内容。但是，该方式操作繁琐，且在后续需要调换或增删约束时很容易漏改某些tag，导致子公式编号混乱，而又不容易察觉。

本模板提供的\shadcmd{subeqtag[<子公式编号标签>]}命令彻底避免了上述问题。使用者只需要在对应的约束后插入\shadcmd{subeqtag}，即可赋予该约束与当前主公式编号保持一致的次级编号。而且，对连续多个约束使用该命令会\textbf{自动生成}递增的子公式编号，交换约束顺序编号也会自行更新，断不会出错。

如果需要在正文中引用某个子公式编号，那么可以像往常一样在\shadcmd{subeqtag}之后使用\shadcmd{label\{<编号标签>\}}，或者直接指定\shadcmd{subeqtag[<子公式编号标签>]}的可选参数，非常人性化。

下面的源码将产生式\eqref{eq: obj 1}\textasciitilde \eqref{eq: constriant gamma}对应的例子，其中式\eqref{eq: constraint x}和\eqref{eq: constriant gamma}使用了\shadcmd{subeqtag[<子公式编号标签>]}的可选参数。


\begin{verbatim}
    \begin{align}
        \max \log \left(x^2 + y^2 + z^2 + v^2 + g^2 + m^2 + k^2\right)
        \label{eq: obj 1} \\
        \text{s.t.} \quad x \leq 1, \subeqtag[eq: constraint x] \\
        y \leq 2, \subeqtag \\
        z \leq 4, \subeqtag
    \end{align}
    
    \begin{align}
        \min \left(\boldsymbol{\alpha} + \beta + \gamma_{中文}\right)^2
        \label{eq: obj 2} \\
        \text{s.t.} \quad \boldsymbol{\alpha} \leq 9, \subeqtag \\
        \beta \geq -10, \subeqtag \\
        \gamma_{中文} \geq 8, \subeqtag[eq: constriant gamma]
    \end{align}
\end{verbatim}
\begin{align}
    \thinmuskip=-3mu \medmuskip=-2mu \thickmuskip=-1mu
    \max \log \left(x^2 + y^2 + z^2 + v^2 + g^2 + m^2 + k^2\right) \label{eq: obj 1} \\
    \text{s.t.} \quad x \leq 1, \subeqtag[eq: constraint x] \\
    y \leq 2, \subeqtag \\
    z \leq 4, \subeqtag
\end{align}
\begin{align}
    \min \left(\boldsymbol{\alpha} + \beta + \gamma_{中文}\right)^2 \label{eq: obj 2} \\
    \text{s.t.} \quad \boldsymbol{\alpha} \leq 9, \subeqtag \\
    \beta \geq -10, \subeqtag \\
    \gamma_{中文} \geq 8, \subeqtag[eq: constriant gamma]
\end{align}


$\boldsymbol{x}^2 + \mathrm{y}^2 + z^2 + \mathcal{B}^2 + \mathcal{M}^2 + \mathbb{I}^2 + v^2 + g^2 + m^2 + k^2$

$\check{\boldsymbol{C}}^t_n + \hat{p}^j_{n, s} + \tilde{r}^t_{u,b} + \overline{i}^t_{x, y} + \acute{\alpha}^t_v + \acute{\boldsymbol{\alpha}}^t + \hat{\boldsymbol{\alpha}}^t + f^t_j + f^{"t}_j + l^t_j + l^{"t}_j$

$\xlongequal{uvw}, \cdot, \cdots, \xLeftrightarrow{xyz}, \mathcal{N} \triangleq \left\{1, \dots, N\right\}, \times, \partial, \emptyset, \in, \subseteq, \leftarrow, \rightarrow, \leq, \geq$

有两点需要提醒：
\begin{itemize}
    \item \shadcmd{subeqtag[<子公式编号标签>]}的可选参数全文不可重复定义，因为它本质上还是调用的\shadcmd{label\{<编号标签>\}}。
    \item 尽管使用\shadcmd{subeqtag[<子公式编号标签>]}的可选参数指定的标签本质上是基于\shadcmd{label\{<编号标签>\}}进行的封装，TeXstudio编辑器在使用\shadcmd{ref\{<编号标签>\}}或\shadcmd{eqref\{<编号标签>\}}却不会自动弹出这些标签的选项，需要手动输入；而如果是直接用\shadcmd{label\{<编号标签>\}}指定的标签，引用时会出现在选项提示中，可以直接选择。这是\shadcmd{subeqtag[<子公式编号标签>]}在接受可选参数后的不便之处，可惜我并不知道该如何解决。
\end{itemize}

	
\section{排版及微调数学公式}

对于不熟悉数学符号与LaTeX源码对应关系的用户，请参考\href{https://www.cnblogs.com/1024th/p/11623258.html}{\color{DarkRed}\CJKunderline*[thickness=0.5bp, format=\color{DarkRed}]{LaTeX公式手册（全网最全）@樱花赞}}，此处不再赘述。

当出现包含公式较多的数学模型或行间公式组，且当页剩余的排版空间无法完整容纳它们时，用户可以在导言区或文档开头使用\shadcmd{allowdisplaybreaks[<跨页倾向值>]}来允许跨页排版行间公式组，从而避免该页底部出现大幅空白。该命令的可选参数可取的值有\shad{1,2,3,4}，值越大表示跨页排版的倾向越高。

当某些行间公式过长以致于超出页面边界时，用户可以在\shad{equation}环境中嵌套\shad{aligned}环境将之调整为多行排版。此外，若公式超出页面边界的部分不多，也可以在\shad{equation}环境内通过调整公式中数学符号间的三种间距\shadcmd{thinmuskip}、\shadcmd{medmuskip}、\shadcmd{thickmuskip}来略微压缩公式的排版长度。这三项长度的单位只能是\shad{mu}，用户可对比公式\eqref{eq: 超长行间公式}和\eqref{eq: 微调公式中数学符号间距}的排版效果：两者的内容完全相同，但式\eqref{eq: 微调公式中数学符号间距}调整了上述三项长度。
\begin{equation} \label{eq: 超长行间公式}
    \sum_{k=a}^{b} f(k) = \int_{a}^{b} f(x) \, dx + \frac{f(a) + f(b)}{2} + \sum_{m=1}^{\lfloor p/2 \rfloor} \frac{B_{2m}}{(2m)!} \left( f^{(2m-1)}(b) - f^{(2m-1)}(a) \right) + R_p,
\end{equation}
\begin{equation} \label{eq: 微调公式中数学符号间距}
    \thinmuskip=-1mu \medmuskip=-1mu \thickmuskip=0mu
    \sum_{k=a}^{b} f(k) = \int_{a}^{b} f(x) \, dx + \frac{f(a) + f(b)}{2} + \sum_{m=1}^{\lfloor p/2 \rfloor} \frac{B_{2m}}{(2m)!} \left( f^{(2m-1)}(b) - f^{(2m-1)}(a) \right) + R_p,
\end{equation}


\section{在标题中排版数学符号\texorpdfstring{$\tilde{r}^t_{u,b}, \acute{\alpha}^t_v, \check{\boldsymbol{C}}^t_n$}{示例}}

尽管我不建议在标题中排版数学符号（因为规范甚至不建议在标题中排版英文缩略词），但如果非排版不可，那可参考本节标题的做法，使用\href{https://mirrors.tuna.tsinghua.edu.cn/CTAN/macros/latex/contrib/hyperref/doc/hyperref-doc.pdf}{\ttfamily\color{DarkRed}\CJKunderline*[thickness=0.5bp, format=\color{DarkRed}]{hyperref}}宏包（模板已载入该宏包）提供的\shadcmd{texorpdfstring\{<TeXstring>\}\{<PDFstring>\}}命令，该命令的具体用法参考这个帖子：\href{https://blog.csdn.net/qq_42679415/article/details/139592054}{\ttfamily\color{DarkRed}\CJKunderline*[thickness=0.5bp, format=\color{DarkRed}]{texorpdfstring使用方法}}。


\section{排版化学方程式}

本模板基于\href{https://mirrors.cloud.tencent.com/CTAN/macros/latex/contrib/mhchem/mhchem.pdf}{\ttfamily\color{DarkRed}\CJKunderline*[thickness=0.5bp, format=\color{DarkRed}]{mhchem}}和\href{https://mirrors.bfsu.edu.cn/CTAN/macros/generic/chemfig/chemfig-en.pdf}{\ttfamily\color{DarkRed}\CJKunderline*[thickness=0.5bp, format=\color{DarkRed}]{chemfig}}宏包来排版化学方程式、结构式和键线式等。相关命令的使用可参考宏包的官方文档或\href{https://www.luogu.com.cn/user/45443}{\color{DarkRed}\CJKunderline*[thickness=0.5bp, format=\color{DarkRed}]{@codesonic}}的文章：\href{https://www.luogu.com.cn/article/o7mlv3w8}{\color{DarkRed}\CJKunderline*[thickness=0.5bp, format=\color{DarkRed}]{用LaTeX写化学方程式}}。下方示例取自该文章，感谢。

\ce{2H2 + O2 ->T[点燃] 2H2O}，\quad \ce{N2 + 3H2 <=>T[高温、加压][催化剂] 2NH3}

\ce{SO4^2- + Ba^2+ -> BaSO4 v}，\quad \ce{2HCl + Na2CO3 -> H2O + CO2 ^ + 2NaCl}

\ce{^{227}_{90}Th+}，\quad \ce{KCr(SO4)2 * 12H2O}，\quad \ce{C6H5-CHO}，\quad \ce{X=Y#Z}

\ce{\chemfig{CH_3C(=[1]O)(-[7]CH_3)} + \chemfig{CN(-[2]H)} ->T[催化剂] \chemfig{CH_3(-[0]C(-[2]OH)(-[0]CN)(-[6]CH3))}}

\ce{\chemfig{CH_3C(=[1,0.7]O)(-[7,0.7]CH_3)} + \chemfig{CN(-[2,0.7]H)} ->T[催化剂] \chemfig{CH_3(-[0,0.7]C(-[2,0.7]OH)(-[0,0.7]CN)(-[6,0.7]CH3))}}

\chemfig{[,0.7]*6(-=(*5(-N(-H)-=(-[:30]CH_2CH_2NHCOCH_3)--))-=-(-H_3CO)=)}

\chemfig{[,0.7]*6(-=(*5(-[0.7]N(-H)-=(-[:30,0.7](-[:330,0.7](-[:30,0.7]N(-[2,0.7]H)(-[:330,0.7](=[6,0.7]O)(-[:30,0.7])))))--))-=-(-[,0.7]O(-[1,0.7]))=)}


\section{引用}

对“公式”、“图片”、“表格”、“伪码”、“定义”、“公理”、“定理”、“命题”、“推论”、“引理”、“示例”、“假设”、“证明”等编号的引用直接用\shadcmd{ref\{<编号label>\}}即可，其中需要带括号公式编号则使用\shadcmd{eqref\{<公式label>\}}。

若要对子图题编号进行完整引用直接使用\shadcmd{ref\{<子图题标签>\}}即可，\shad{DissertUESTC}文档类默认生成形如\textcolor{DodgerBlue}{1-1(a)}的完整编号，但若指定了文档类的\shad{subfigsimple}选项，则会生成形如\textcolor{DodgerBlue}{1-1a}的完整编号（\textit{注：学位论文撰写规范中并未明确说明引用子图编号应该采用哪种形式，但我翻了本中文专著，里面采用了\textcolor{DodgerBlue}{1-1(a)}的形式，故而将之设为默认样式}）；反之，若只希望单独引用子图题编号，比如在图题结尾按编号添加子图题文本，则需要使用\shadcmd{subref\{<子图题标签>\}}，它将生成形如\textcolor{DodgerBlue}{(a)}的单独编号。

对参考文献的行内引用直接使用\shadcmd{cite\{<参考文献label>\}}，以上标形式引用则使用\shadcmd{citess\{<参考文献label>\}}。






